The spacecraft keeps a fixed orientation and is a rigid body, so all its
points have the same angular velocity $\omega$. The orbital equation gives
\[
K\frac{m_{NC}M}{r^2}=m_{NC}\omega^2r,\qquad \omega^2=\frac{KM}{r^3}.
\]
The balls' centripetal accelerations are
\[
a_{cp(a)}=a_{cp(d)}=\omega^2r,\quad
a_{cp(b)}=\omega^2(r-d),\quad a_{cp(c)}=\omega^2(r+d).
\]
For (a), gravity alone supplies the centripetal force,
\[
T_{(a)}=K\frac{mM}{r^2}-m\omega^2r=0.
\]
For (b), use $d\ll r$ and $(1-d/r)^{-2}\approx1+2d/r$,
\[
F_{ag(b)}=K\frac{mM}{(r-d)^2}
\approx K\frac{mM}{r^2}\left(1+\frac{2d}{r}\right),
\]
\[
T_{(b)}=F_{ag(b)}-m\omega^2(r-d)
=3K\frac{mMd}{r^3}>0.
\]
For (c), $(1+d/r)^{-2}\approx1-2d/r$,
\[
F_{ag(c)}\approx K\frac{mM}{r^2}\left(1-\frac{2d}{r}\right),\qquad
T_{(c)}=m\omega^2(r+d)-F_{ag(c)}=3K\frac{mMd}{r^3}>0.
\]
% sic: An intermediate line in the official solution prints T_(c)=-3KmMd/r^3, immediately followed by the positive result reproduced here.
For (d),
\[T_{(d)}=m\omega^2r-K\frac{mM}{r^2}=0.\]
Thus (a) and (d) are weightless with untensioned wires, while (b) and (c)
have tensioned wires. The answer is (a), with $K=G$.
